% (c) Nikita Lisitsa, lisyarus@gmail.com, 2026

\documentclass[10pt]{beamer}

\usepackage[T2A]{fontenc}
\usepackage[russian]{babel}
\usepackage{minted}

\usepackage[most]{tcolorbox}
\tcbuselibrary{minted}

\usepackage{graphicx}
\graphicspath{ {./images/} }

\usepackage{adjustbox}

\usepackage{color}
\usepackage{soul}
\usepackage{multicol}
\usepackage{hyperref}
\usepackage{tikz}
\usetikzlibrary{arrows.meta,positioning,calc}

\usetheme{metropolis}

\definecolor{red}{rgb}{1,0,0}
\definecolor{codebg}{RGB}{29,35,49}
\setminted{fontsize=\footnotesize}

\makeatletter
\newcommand{\slideimage}[1]{
  \begin{figure}
    \begin{adjustbox}{width=\textwidth, totalheight=\textheight-2\baselineskip-2\baselineskip,keepaspectratio}
      \includegraphics{#1}
    \end{adjustbox}
  \end{figure}
}
\makeatother

\title{Язык программирования C++}
\subtitle{Лекция 19: Инициализация объектов, final, tuple, regex, анонимные функции, function, type erasure, многопоточное программирование}
\date{2026}
\author{Никита Лисица}

\setbeamertemplate{footline}[frame number]

\begin{document}

\frame{\titlepage}

% TODO:
% + инициализация объектов
% + final
% + tuple
% + лямбды
% + function, type erasure
% * многопоточное программирование

\begin{frame}[fragile]
\frametitle{Инициализация объектов}
\begin{itemize}
\usemintedstyle{tango}
\item Мы знаем, что основной способ инициализировать объект в C++ -- вызвать конструктор
\pause
\item На самом деле, всё несколько сложнее :)
\pause
\item Default initialization: \mintinline{cpp}|Type x;|
\pause
\item Value initialization: \mintinline{cpp}|Type x{};|
\pause
\item Direct initialization: \mintinline{cpp}|Type x(a, b);|
\pause
\item Copy initialization: \mintinline{cpp}|Type x = a;|
\pause
\item List initialization: \mintinline{cpp}|Type x{a, b};| или \mintinline{cpp}|Type x = {a, b};|
\end{itemize}
\end{frame}

\begin{frame}[fragile]
\frametitle{Default vs value initialization}
\begin{itemize}
\usemintedstyle{tango}
\item Default initialization \mintinline{cpp}|Type x;| для классов вызывает конструктор без аргументов, для примитивных типов ничего не делает
\pause
\item \textbf{NB}: если у класса нет конструктора, компилятор сам создаст конструктор, рекурсивно вызывающий конструкторы всех полей (но для примитивных типов они тоже ничего не делают)
\pause
\item Value initialization \mintinline{cpp}|Type x{};| зануляет объект, потом для классов вызывает конструктор без аргументов
\end{itemize}
\end{frame}

\begin{frame}[fragile]
\frametitle{Default vs value initialization}
\begin{itemize}
\usemintedstyle{lightbulb}
\begin{tcblisting}{listing engine=minted, minted language=cpp, minted options={fontsize=\scriptsize}, listing only, colback=codebg, colframe=codebg, rounded corners}
// Нет конструктора без аргументов
// Компилятор сам создаст конструктор,
// который ничего не будет делать
// (trivial constructor)
struct point {
  int x, y;
};

// Default initialization
// Значение остаётся не проинициализированным
int i;
// Поля остаются не проинициализированы
point p;

// Value initialization
// Значение будет нулём
int i{};
// Значение будет нулём
point p{};
\end{tcblisting}
\end{itemize}
\end{frame}

\begin{frame}[fragile]
\frametitle{List initialization}
\begin{itemize}
\usemintedstyle{tango}
\item Синтаксис \mintinline{cpp}|Type x = {a, b, c};| может сделать одну из нескольких вещей:
\pause
\item Вызвать конструктор \mintinline{cpp}|Type| от аргументов \mintinline{cpp}|(a, b, c)|
\pause
\item Если у \mintinline{cpp}|Type| нет конструктора, проициниализировать все поля значениями из \mintinline{cpp}|{a, b, c}| по порядку
\pause
\item Если \mintinline{cpp}|Type| это массив, проинициализировать его элементы значениями из \mintinline{cpp}|{a, b, c}| по порядку
\pause
\item Вызвать конструктор от \mintinline{cpp}|std::initializer_list|
\end{itemize}
\end{frame}

\begin{frame}[fragile]
\frametitle{List initialization: примеры}
\begin{itemize}
\usemintedstyle{lightbulb}
\begin{tcblisting}{listing engine=minted, minted language=cpp, minted options={fontsize=\scriptsize}, listing only, colback=codebg, colframe=codebg, rounded corners}
struct point {
  int x, y;
  point(int x, int y)
    : x(x), y(y)
  {}
};

// Вызов конструктора point::point(int, int)
point p = {1, 2};
\end{tcblisting}
\pause
\usemintedstyle{lightbulb}
\begin{tcblisting}{listing engine=minted, minted language=cpp, minted options={fontsize=\scriptsize}, listing only, colback=codebg, colframe=codebg, rounded corners}
struct point {
  int x, y;
};

// Aggregate initialization
// Эквивалентно p.x = 1, p.y = 2
point p = {1, 2};
\end{tcblisting}
\end{itemize}
\end{frame}

\begin{frame}[fragile]
\frametitle{List initialization: примеры}
\begin{itemize}
\usemintedstyle{lightbulb}
\item При aggregate initialization можно указывать явно имена полей (\textit{designated initializers})
\begin{tcblisting}{listing engine=minted, minted language=cpp, minted options={fontsize=\scriptsize}, listing only, colback=codebg, colframe=codebg, rounded corners}
struct point {
  int x, y;
};

point p = {.x = 1, .y = 2};
\end{tcblisting}
\pause
\usemintedstyle{lightbulb}
\begin{tcblisting}{listing engine=minted, minted language=cpp, minted options={fontsize=\scriptsize}, listing only, colback=codebg, colframe=codebg, rounded corners}
// Инициализация массива
int array[3] = {1, 2, 3};
\end{tcblisting}
\end{itemize}
\end{frame}

\begin{frame}[fragile]
\frametitle{std::initializer\_list}
\begin{itemize}
\usemintedstyle{tango}
\item \mintinline{cpp}|std::initializer_list<T>| -- специальный шаблонный класс в стандартной библиотеке для инициализации контейнеров
\pause
\item Обычно ссылается на массив элементов (\mintinline{cpp}|T* data, size_t size|), где-то созданных компилятором
\pause
\usemintedstyle{lightbulb}
\begin{tcblisting}{listing engine=minted, minted language=cpp, minted options={fontsize=\scriptsize}, listing only, colback=codebg, colframe=codebg, rounded corners}
// Вызывает конструктор от
// std::initializer_list<int>
std::vector<int> vec = {1, 2, 3, 4};
\end{tcblisting}
\pause
\item Все стандартные контейнеры имеют такие конструкторы
\end{itemize}
\end{frame}

\begin{frame}[fragile]
\frametitle{std::initializer\_list}
\begin{itemize}
\usemintedstyle{tango}
\item Можно создать такой конструктор и своему классу:
\usemintedstyle{lightbulb}
\begin{tcblisting}{listing engine=minted, minted language=cpp, minted options={fontsize=\scriptsize}, listing only, colback=codebg, colframe=codebg, rounded corners}
class my_container {
public:
  my_container(std::initializer_list<int> l) {
    // l.begin(), l.end() -- указатели на начало/конец
    // l.data(), l.size(), l.empty() -- как у std::vector
  }
};

my_container c = {1, 2, 3, 4};
\end{tcblisting}
\end{itemize}
\end{frame}

\begin{frame}[fragile]
\frametitle{final}
\begin{itemize}
\usemintedstyle{tango}
\item Можно запретить наследоваться от класса:
\usemintedstyle{lightbulb}
\begin{tcblisting}{listing engine=minted, minted language=cpp, minted options={fontsize=\scriptsize}, listing only, colback=codebg, colframe=codebg, rounded corners}
class A final {
  // ...
};

// Ошибка компиляции: от A нельзя
// наследоваться
class B : public A {};
\end{tcblisting}
\end{itemize}
\end{frame}

\begin{frame}[fragile]
\frametitle{final}
\begin{itemize}
\usemintedstyle{tango}
\item Можно запретить переопределять виртуальный метод:
\usemintedstyle{lightbulb}
\begin{tcblisting}{listing engine=minted, minted language=cpp, minted options={fontsize=\scriptsize}, listing only, colback=codebg, colframe=codebg, rounded corners}
class A {
public:
  virtual void foo();
};

class B : public A {
public:
  void foo() final;
};

// Ошибка компиляции: foo() нельзя
// переопределять
class C : public B {
public:
  void foo() override;
};
\end{tcblisting}
\end{itemize}
\end{frame}

\begin{frame}[fragile]
\frametitle{final}
\begin{itemize}
\usemintedstyle{tango}
\item Зачем использовать \mintinline{cpp}|final|? \pause Для оптимизации и для защиты от ошибок
\usemintedstyle{lightbulb}
\begin{tcblisting}{listing engine=minted, minted language=cpp, minted options={fontsize=\scriptsize}, listing only, colback=codebg, colframe=codebg, rounded corners}
class A {
public:
  virtual void foo();
};

class B : public A {
public:
  void foo() final;
};

void call_foo(B* ptr) {
  // Можно не использовать dynamic dispatch:
  // мы знаем, что это именно B::foo(), так
  // как её нельзя переопределить
  ptr->foo();
}
\end{tcblisting}
\end{itemize}
\end{frame}

\begin{frame}[fragile]
\frametitle{std::tuple}
\begin{itemize}
\usemintedstyle{tango}
\item \mintinline{cpp}|std::tuple| -- реализация \textit{кортежа}, т.е. упорядоченного набора произвольных элементов
\pause
\item \textbf{NB}: во многих языках tuple -- встроенная конструкция, но в C++ она реализована в стандартной библиотеке
\pause
\item Это шаблонный класс, параметризующийся хранимыми типами, например \mintinline{cpp}|std::tuple<int, float, std::string>|
\pause
\usemintedstyle{lightbulb}
\begin{tcblisting}{listing engine=minted, minted language=cpp, minted options={fontsize=\scriptsize}, listing only, colback=codebg, colframe=codebg, rounded corners}
std::tuple<int, float> tup1{15, 3.14f};

// Используем вывод шаблонных параметров классов
// -> std::tuple<char, double>
std::tuple tup2{'a', 10.0};

// Получить значение по compile-time индексу:
std::get<0>(tup1); // == 15
std::get<1>(tup2); // == 10.0
\end{tcblisting}
\end{itemize}
\end{frame}

\begin{frame}[fragile]
\frametitle{std::tuple}
\begin{itemize}
\usemintedstyle{tango}
\item Основное применение \mintinline{cpp}|std::tuple| -- сохранить передать неизвестный набор значений куда-либо, обычно в шаблонном коде
\pause
\item Лучше не использовать \mintinline{cpp}|std::tuple| в нормальных интерфейсах:
\usemintedstyle{lightbulb}
\begin{tcblisting}{listing engine=minted, minted language=cpp, minted options={fontsize=\scriptsize}, listing only, colback=codebg, colframe=codebg, rounded corners}
// Что такое int и float?
std::tuple<int, float> get_time();

// Лучше иметь структуры с явными
// именами полей:
struct time {
  int seconds;
  float milliseconds;
};

time get_time();
\end{tcblisting}
\end{itemize}
\end{frame}

\begin{frame}[fragile]
\frametitle{Анонимные функции}
\begin{itemize}
\usemintedstyle{tango}
\item При работе со стандартной библиотекой часто приходится написать функцию/функциональный объект: компаратор для \mintinline{cpp}|std::sort| или \mintinline{cpp}|std::map|, предикат для \mintinline{cpp}|std::find_if|, и т.п.
\pause
\usemintedstyle{lightbulb}
\begin{tcblisting}{listing engine=minted, minted language=cpp, minted options={fontsize=\scriptsize}, listing only, colback=codebg, colframe=codebg, rounded corners}
struct less_by_first_character
{
  bool operator()(
    const std::string& s1,
    const std::string& s2) {
    return s1[0] < s2[0];
  }
};

std::sort(lines.begin(), lines.end(),
  less_by_first_character());
\end{tcblisting}
\pause
\item Писать целый класс неудобно: много лишнего кода (т.н. \textit{boilerplate}), логика `размазана' по коду
\end{itemize}
\end{frame}

\begin{frame}[fragile]
\frametitle{Анонимные функции}
\begin{itemize}
\usemintedstyle{tango}
\item Решение: \textit{анонимные функции}, они же -- \textit{лямбда-функции} или просто \textit{лямбды}
\pause
\item Концепция пришла из \textit{лямбда-исчисления} (Алонзо Чёрч, 1930-е) -- базис для функционального программирования
\pause
\item \( Y = \lambda f.(\lambda x.f(x\ x))\ (\lambda x.f(x\ x)) \)
\pause
\item В императивных языках лямбда-функции -- синтаксический сахар, т.е. удобный способ выразить часто повторяющуюся конструкцию
\end{itemize}
\end{frame}

\begin{frame}[fragile]
\frametitle{Анонимные функции}
\begin{itemize}
\usemintedstyle{tango}
\item Общий синтаксис: \mintinline{cpp}|[список захвата](аргументы){ тело функции; }|
\pause
\usemintedstyle{lightbulb}
\begin{tcblisting}{listing engine=minted, minted language=cpp, minted options={fontsize=\scriptsize}, listing only, colback=codebg, colframe=codebg, rounded corners}
auto sqr = [](int x){ return x * x; };
auto sum = [](int x, int y){ return x + y; };
\end{tcblisting}
\pause
\item Компилятор превратит лямбды в функциональные объекты:
\usemintedstyle{lightbulb}
\begin{tcblisting}{listing engine=minted, minted language=cpp, minted options={fontsize=\scriptsize}, listing only, colback=codebg, colframe=codebg, rounded corners}
struct sqr_t {
  int operator()(int x) const { return x * x; };
};

struct sum_t {
  int operator()(int x, int y) const { return x + y; };
};

auto sqr = sqr_t{};
auto sum = sum_t{};
\end{tcblisting}
\end{itemize}
\end{frame}

\begin{frame}[fragile]
\frametitle{Анонимные функции}
\begin{itemize}
\usemintedstyle{tango}
\item Их удобно использовать как компараторы/предикаты:
\usemintedstyle{lightbulb}
\begin{tcblisting}{listing engine=minted, minted language=cpp, minted options={fontsize=\scriptsize}, listing only, colback=codebg, colframe=codebg, rounded corners}
// Отсортировать по возрастанию модуля
std::sort(array.begin(), array.end(),
  [](int x, int y){ abs(x) < abs(y); });

// Найти элемент, кратный 8
std::find_if(array.begin(), array.end(),
  [](int x){ return (x % 8) == 0; });
\end{tcblisting}
\end{itemize}
\end{frame}

\begin{frame}[fragile]
\frametitle{Capture list}
\begin{itemize}
\usemintedstyle{tango}
\item Удобство функциональных объектов в том, что они могут \textit{хранить данные} и использовать их в своём \mintinline{cpp}|operator()|:
\usemintedstyle{lightbulb}
\begin{tcblisting}{listing engine=minted, minted language=cpp, minted options={fontsize=\scriptsize}, listing only, colback=codebg, colframe=codebg, rounded corners}
struct less_than_ref {
  int ref;
  bool operator()(int x) const { return x < ref; }
};

std::partition(array.begin(), array.end(),
  less_than_ref{15});
\end{tcblisting}
\pause
\item У лямбд для этого есть \textit{список захвата (capture list)}:
\usemintedstyle{lightbulb}
\begin{tcblisting}{listing engine=minted, minted language=cpp, minted options={fontsize=\scriptsize}, listing only, colback=codebg, colframe=codebg, rounded corners}
int ref = 15;
std::partition(array.begin(), array.end(),
  [ref](int x){ return x < ref; });
\end{tcblisting}
\end{itemize}
\end{frame}

\begin{frame}[fragile]
\frametitle{Capture list}
\begin{itemize}
\usemintedstyle{tango}
\item `Захватить' переменную в список захвата можно как по значению, так и по ссылке
\pause
\item Захват по значению \mintinline{cpp}|[x](...){...}| создаёт копию переменной внутри лямбды
\pause
\item Захват по ссылке \mintinline{cpp}|[&x](...){...}| создаёт ссылку на переменную внутри лямбды
\pause
\item \textbf{NB}: в захвате по ссылке выражение \mintinline{cpp}|&x| не означает указатель на переменную!
\pause
\usemintedstyle{lightbulb}
\begin{tcblisting}{listing engine=minted, minted language=cpp, minted options={fontsize=\scriptsize}, listing only, colback=codebg, colframe=codebg, rounded corners}
int count = 0;
int ref = 15;
std::partition(begin, end, [&count, ref](int x){
  ++count;
  return x < ref;
});
\end{tcblisting}
\end{itemize}
\end{frame}

\begin{frame}[fragile]
\frametitle{Capture list}
\begin{itemize}
\usemintedstyle{tango}
\item Эквивалентный код без лямбд:
\usemintedstyle{lightbulb}
\begin{tcblisting}{listing engine=minted, minted language=cpp, minted options={fontsize=\scriptsize}, listing only, colback=codebg, colframe=codebg, rounded corners}
struct lambda_t {
  int& count;
  int ref;

  bool operator()(int x) const {
    ++count;
    return x < ref;
  }
};

int count = 0;
int ref = 15;
std::partition(begin, end, lambda_t{count, ref});
\end{tcblisting}
\end{itemize}
\end{frame}

\begin{frame}[fragile]
\frametitle{Capture list}
\begin{itemize}
\usemintedstyle{tango}
\item Элементы списка захвата не обязаны быть существующими переменными, но тогда их надо руками проинициализировать:
\usemintedstyle{lightbulb}
\begin{tcblisting}{listing engine=minted, minted language=cpp, minted options={fontsize=\scriptsize}, listing only, colback=codebg, colframe=codebg, rounded corners}
int ref = 10;
auto func = [y = ref * 2](int x){ return x < y; };
\end{tcblisting}
\pause
\usemintedstyle{tango}
\item По умолчанию, сгенерированный \mintinline{cpp}|operator()| помечен как \mintinline{cpp}|const| и не может менять значения в списке захвата
\pause
\item Это можно изменить ключевым словом \mintinline{cpp}|mutable|:
\usemintedstyle{lightbulb}
\begin{tcblisting}{listing engine=minted, minted language=cpp, minted options={fontsize=\scriptsize}, listing only, colback=codebg, colframe=codebg, rounded corners}
auto func = [count = 0](int x) mutable {
  ++count;
  return x < 15;
};
\end{tcblisting}
\end{itemize}
\end{frame}

\begin{frame}[fragile]
\frametitle{Capture list}
\begin{itemize}
\usemintedstyle{tango}
\item Можно захватить всё, что используется в лямбде, автоматически по ссылке \mintinline{cpp}|[&]| или по значению \mintinline{cpp}|[=]|:
\usemintedstyle{lightbulb}
\begin{tcblisting}{listing engine=minted, minted language=cpp, minted options={fontsize=\scriptsize}, listing only, colback=codebg, colframe=codebg, rounded corners}
int ref = 10;

// Автоматический захват по значению
auto func1 = [=](int x){ return x < ref; };

// Автоматический захват по ссылке
auto func2 = [&](int x){ return x < ref; };
\end{tcblisting}
\pause
\item \textbf{NB}: Лучше не пользоваться автоматическим захватом, тяжело отследить битые ссылки
\end{itemize}
\end{frame}

\begin{frame}[fragile]
\frametitle{Анонимные функции: возвращаемое значение}
\begin{itemize}
\usemintedstyle{tango}
\item Тип возвращаемого значения у лямбд выводится автоматически
\pause
\item Можно указать его явно, чтобы избежать ошибок или неявных преобразований:
\usemintedstyle{lightbulb}
\begin{tcblisting}{listing engine=minted, minted language=cpp, minted options={fontsize=\scriptsize}, listing only, colback=codebg, colframe=codebg, rounded corners}
auto sqr = [](int x) -> int {
  return x * x;
};

auto cmp = [](int x, int y) -> bool {
  return abs(x) < abs(y);
};
\end{tcblisting}
\end{itemize}
\end{frame}

\begin{frame}[fragile]
\frametitle{Анонимные функции: применение}
\begin{itemize}
\usemintedstyle{tango}
\item Лямбды можно использовать не только как предикаты или компараторы
\pause
\item Настройка поведения сложного шаблонного кода
\pause
\item Callback'и (обработчики событий) в асинхронном коде (сеть, UI)
\pause
\item Библиотека \nolinkurl{Boost.Hana} реализует метапрограммрование на лямбдах
\end{itemize}
\end{frame}

\begin{frame}[fragile]
\frametitle{std::function}
\begin{itemize}
\usemintedstyle{tango}
\item Если мы хотим куда-то принимать функцию или лямбду, у нас есть два варианта:
\pause
\begin{itemize}
\item Принимать указатель на функцию -- тогда нельзя передать лямбду
\pause
\item Сделать код шаблонным -- тогда его придётся положить в заголовочный файл, это замедлит компиляцию, и его нельзя положить в динамическую библиотеку
\end{itemize}
\pause
\item Решение: \mintinline{cpp}|std::function|
\end{itemize}
\end{frame}

\begin{frame}[fragile]
\frametitle{std::function}
\begin{itemize}
\usemintedstyle{tango}
\item \mintinline{cpp}|std::function<R(Args...)>| -- специальный шаблонный класс, позволяющий хранить любой объект, который можно вызвать от списка аргументов \mintinline{cpp}|(Args...)|, и который возвращает тип \mintinline{cpp}|R|
\pause
\usemintedstyle{lightbulb}
\begin{tcblisting}{listing engine=minted, minted language=cpp, minted options={fontsize=\scriptsize}, listing only, colback=codebg, colframe=codebg, rounded corners}
std::function<int(int)> sqr = [](int x){
  return x * x;
};

std::function<int(int, int)> sum = [](int x, int y){
  return x + y;
};
\end{tcblisting}
\end{itemize}
\end{frame}

\begin{frame}[fragile]
\frametitle{std::function}
\begin{itemize}
\usemintedstyle{tango}
\item \mintinline{cpp}|std::function| с конкретными параметрами -- конкретный тип, его можно использовать в нешаблонных интерфейсах классов и библиотек
\pause
\usemintedstyle{lightbulb}
\begin{tcblisting}{listing engine=minted, minted language=cpp, minted options={fontsize=\scriptsize}, listing only, colback=codebg, colframe=codebg, rounded corners}
class UIController {
public:
  virtual onKeyDown(std::function<void(char)> callback) = 0;
  virtual onKeyUp  (std::function<void(char)> callback) = 0;

  virtual onMouseMove
    (std::function<void(int, int)> callback) = 0;

  virtual ~UIController() {}
};

controller->onKeyDown([](char key){
  if (key == 'W') {
    moveForward();
  }
});
\end{tcblisting}
\end{itemize}
\end{frame}

\begin{frame}[fragile]
\frametitle{std::function: реализация}
\begin{itemize}
\usemintedstyle{lightbulb}
\begin{tcblisting}{listing engine=minted, minted language=cpp, minted options={fontsize=\scriptsize}, listing only, colback=codebg, colframe=codebg, rounded corners}
template <typename R, typename ... Args>
struct function_base {
  virtual R call(Args...) = 0;
};

template <typename R, typename ... Args, typename Object>
struct function_impl : base<R, Args...> {
  Object object;
  R call(Args... args) { return object(args...); }
};

template <typename R, typename ... Args>
struct function {
  template <typename Object>
  function(Object object)
    : impl(new function_impl<R, Args..., Object>(object))
  {}

  R operator()(Args... args) { return impl->call(args); }

private:
  function_base<R, Args...>* impl;
};
\end{tcblisting}
\end{itemize}
\end{frame}

\begin{frame}[fragile]
\frametitle{Type erasure}
\begin{itemize}
\usemintedstyle{tango}
\item \mintinline{cpp}|std::function| -- обёртка для любого объекта, обладающего некоторым интерфейсом (\mintinline{cpp}|operator()| с нужными типами)
\pause
\item Реализация имспользует настоящий интерфейс с виртуальными функциями
\pause
\item Оборачивание конкретного объекта создаёт реализацию этого интерфейса, параметризующуюся типом хранимого объекта
\pause
\item Такой приём называется \textit{type erasure} и часто используется чтобы разменять производительность на скорость компиляции и удобство архитектуры
\end{itemize}
\end{frame}

\begin{frame}[fragile]
\frametitle{Процессы vs потоки}
\begin{itemize}
\usemintedstyle{tango}
\item \textit{Процесс} -- отдельно запущенная программа со своим адресным пространством
\pause
\item \textit{Поток} -- независимо выполняющаяся часть программы, использующая её адресное пространство
\pause
\item Разные \textit{процессы} не могут просто так общаться через память (для этого нужны сокеты/файлы/shared memory)
\pause
\item Один \textit{процесс} может иметь несколько \textit{потоков}, и они все имеют общую память
\pause
\item Единица исполнения кода -- это \textit{поток}
\end{itemize}
\end{frame}

\begin{frame}[fragile]
\frametitle{Виды многозадачности}
\begin{itemize}
\usemintedstyle{tango}
\item В общем случае, потоков может быть больше, чем ядер CPU
\pause
\item \textit{Вытесняющая многозадачность} -- ядро операционной системы выделяет кванты времени каждому потоку и переключается между ними на каждом ядре CPU
\pause
\item \textit{Кооперативная многозадачность} -- выполняющиеся потоки сами решают, когда отдать управление другому потоку
\pause
\item В современных ОС реализована \textit{вытесняющая многозадачность} на уровне потоков всех процессов системы
\pause
\item Современные библиотеки работы с сетью (\nolinkurl{node.js, asyncio}) асинхронные и используют \textit{кооперативную многозадачность}, часто с помощью \textit{корутин}
\end{itemize}
\end{frame}

\begin{frame}[fragile]
\frametitle{Потоки}
\begin{itemize}
\usemintedstyle{tango}
\item При создании потока нужно
\pause
\begin{itemize}
\item Выделить место под его стек
\pause
\item Выделить место под \mintinline{cpp}|thread_local| переменные
\pause
\item Проинициализировать информацию о потоке внутри ОС
\end{itemize}
\pause
\item При переключении потоков (т.н. \textit{переключение контекста, context switch}) нужно
\pause
\begin{itemize}
\item Запомнить состояние всех регистров потока
\pause
\item Загрузить состояние всех регистров нового потока
\pause
\item Обновить MMU в CPU для нового процесса
\end{itemize}
\end{itemize}
\end{frame}

\begin{frame}[fragile]
\frametitle{Потоки}
\begin{itemize}
\usemintedstyle{tango}
\item \(\Longrightarrow\) Создание и переключение потоков -- дорогая операция
\pause
\item Перед использованием потоков нужно подумать, не будут ли накладные расходы выше, чем выигрыш от параллельности
\pause
\item Типичное применение:
\pause
\begin{itemize}
\item Распараллеливание алгоритмов (умножение матриц, сортировка)
\pause
\item Real-time приложения (UI в одном потоке, сложные операции в другом)
\pause
\item Сетевые приложения (несколько потоков, каждый обрабатывает свой запрос)
\end{itemize}
\end{itemize}
\end{frame}

\begin{frame}[fragile]
\frametitle{Потоки в C++}
\begin{itemize}
\usemintedstyle{tango}
\item Стандартный класс \mintinline{cpp}|std::thread| реализует отдельный поток выполнения
\pause
\item Конструктор создаёт и сразу запускает поток
\pause
\item Метод \mintinline{cpp}|join()| ждёт, пока поток закончит выполняться
\pause
\usemintedstyle{lightbulb}
\begin{tcblisting}{listing engine=minted, minted language=cpp, minted options={fontsize=\scriptsize}, listing only, colback=codebg, colframe=codebg, rounded corners}
void my_thread_function();

std::thread th(&my_thread_function);

// Можем выполнять другие действия, пока работает поток
...

th.join()
\end{tcblisting}
\pause
\usemintedstyle{tango}
\item \textbf{NB}: с потоком, вызвавшим функцию \mintinline{cpp}|main|, ничего сделать нельзя!
\end{itemize}
\end{frame}

\begin{frame}[fragile]
\frametitle{Потоки в C++}
\begin{itemize}
\usemintedstyle{tango}
\item В конструктор потока можно передать и лямбду
\usemintedstyle{lightbulb}
\begin{tcblisting}{listing engine=minted, minted language=cpp, minted options={fontsize=\scriptsize}, listing only, colback=codebg, colframe=codebg, rounded corners}
// Лямбда без аргументов - можно опустить круглые скобки
std::thread th([]{
  // Тело функции, которая будет исполняться
  // в отдельном потоке
  ...
});
\end{tcblisting}
\pause
\usemintedstyle{tango}
\item В конструктор потока можно передать аргументы для вызываемой функции:
\usemintedstyle{lightbulb}
\begin{tcblisting}{listing engine=minted, minted language=cpp, minted options={fontsize=\scriptsize}, listing only, colback=codebg, colframe=codebg, rounded corners}
void thread_func(int id);

// Создаст поток и вызовет в нём
// thread_func(10)
std::thread th(&thread_func, 10);
\end{tcblisting}
\end{itemize}
\end{frame}

\begin{frame}[fragile]
\frametitle{Потоки в C++}
\begin{itemize}
\usemintedstyle{tango}
\item Пример: параллельное сложение двух векторов
\usemintedstyle{lightbulb}
\begin{tcblisting}{listing engine=minted, minted language=cpp, minted options={fontsize=\scriptsize}, listing only, colback=codebg, colframe=codebg, rounded corners}
void sum(const int* src1, const int* src2,
  int* dst, int count)
{
  for (int i = 0; i < count; ++i)
    dst[i] = src1[i] + src2[i];
}

void parallel_sum(const int* src1, const int* src2,
  int* dst, int count)
{
  // Тут хорошо бы округлить до размера кеш-линии,
  // см. std::hardware_destructive_interference_size
  int half = count / 2;
  std::thread th1(&sum, src1, src2, dst, half);
  std::thread th2(&sum, src1 + half, src2 + half,
    dst + half, count - half);
  th1.join();
  th2.join();
}
\end{tcblisting}
\end{itemize}
\end{frame}

\begin{frame}[fragile]
\frametitle{Потоки в C++}
\begin{itemize}
\usemintedstyle{tango}
\item \textbf{NB}: начиная с C++17 можно такие вещи не реализовывать руками:
\usemintedstyle{lightbulb}
\begin{tcblisting}{listing engine=minted, minted language=cpp, minted options={fontsize=\scriptsize}, listing only, colback=codebg, colframe=codebg, rounded corners}
void parallel_sum(const int* src1, const int* src2,
  int* dst, int count)
{
  // Вместо лямбды можно std::plus<int>{}
  std::transform(std::execution::par, src1, src1 + count,
    src2, dst, [](int x, int y){ return x + y; });
}
\end{tcblisting}
\end{itemize}
\end{frame}

\end{document}